%% kapa: exam=1 round=27 year=2016 period=1 form=A subject=civil_law subject_ko=민법 seq=10 number=10
%% source: Q-Net 감정평가사 기출문제 (cst00309) artlSeq=5169402 file=extracted/1008/5169402/제27회 감정평가사 1차 문제/제27회 감정평가사 1교시문제(A형).pdf
%% conversion: kordoc 4.24.51 -> tools/gen1.py (auto); review=proofread-sonnet; flags=none
\begin{question}{10}
권한을 넘은 표현대리에 관한 설명으로 옳지 않은 것은? (다툼이 있으면 판례에 따름)
\begin{choices}
  \item 복대리인 선임권이 없는 대리인에 의하여 선임된 복대리인의 권한도 기본대리권이 될 수 있다.
  \item 정당한 이유의 유무는 대리행위 당시와 그 이후의 사정을 고려하여 판단한다.
  \item 기본대리권은 표현대리행위와 동종 또는 유사할 필요가 없다.
  \item 권한을 넘은 표현대리는 법정대리에도 적용된다.
  \item 대리행위가 대리권을 제한하는 강행규정을 위반하여 권한을 넘은 경우에는 표현대리가 인정되지 않는다.
\end{choices}
\end{question}
\answer{2}
